% Source fingerprint: 88b089aaadd01ba58e73e9a153026309df6f8ac9e9e35b19f12fcd69aca8859f
% T14: raw TeX cells preserved; no numeric highlighting.
\begin{table}[htbp]
\centering\small\renewcommand{\arraystretch}{1.18}\setlength{\tabcolsep}{4pt}
\caption[容量与维数的三个指数区间]{容量与维数的三个指数区间。前三行讨论本章全空间非齐次容量。最后一行是第26.4节建立的相对容量；两种能量的尺度不同。}
\label{b1:tab:visual-t14}
\begin{tabularx}{\linewidth}{@{}>{\raggedright\arraybackslash}p{3.1cm}>{\raggedright\arraybackslash}p{4.7cm}>{\raggedright\arraybackslash}X@{}}
\toprule
指数 & 小球与单点 & 零测集与零容量集 \\
\midrule
\(1<p<d\) & \(C_p(B_r)\le C(r^d+r^{d-p})\)；单点容量零 & 单点给出非空零容量集；平面片可为零测而正容量 \\
\(p=d>1\) & 需用对数型截断证明单点容量零；上一行幂估计本身不够 & 可数集容量零，但不能把所有 Lebesgue 零测集都忽略 \\
\(p>d\) & 单点容量有统一正下界 & 每个非空集合容量为正；只有空集可为零容量集 \\
相对容量的球对 & \(\operatorname{cap}_p(B_r,B_{2r})\) 与 \(r^{d-p}\) 可比 & 纯梯度能量，不能与全空间非齐次容量混为同一数值 \\
\bottomrule
\end{tabularx}
\end{table}
